\documentclass[11pt]{article}
% Shared preamble for the PNPInverse onboarding docs.
% Each document does % Shared preamble for the PNPInverse onboarding docs.
% Each document does % Shared preamble for the PNPInverse onboarding docs.
% Each document does \input{preamble} after \documentclass.

\usepackage[utf8]{inputenc}
\usepackage[T1]{fontenc}
\usepackage{lmodern}
\usepackage[margin=1in]{geometry}
\usepackage{amsmath, amssymb, amsthm}
\usepackage{mathtools}
\usepackage{empheq}
\usepackage{siunitx}
\usepackage{booktabs}
\usepackage{array}
\usepackage{enumitem}
\usepackage{xcolor}
\usepackage{framed}
\usepackage{microtype}
\usepackage{listings}
\usepackage{tikz}
\usetikzlibrary{arrows.meta, positioning, decorations.pathreplacing, calligraphy}
\usepackage[hidelinks]{hyperref}

% --- branding -------------------------------------------------------------
\definecolor{pnpblue}{RGB}{30,70,120}
\definecolor{pnpgray}{RGB}{90,90,90}
\definecolor{calloutbg}{RGB}{245,245,235}
\definecolor{shadecolor}{RGB}{245,245,235}

% --- callouts (uses framed.sty's shaded environment) ----------------------
\newenvironment{callout}[1]
  {\begin{shaded}\noindent\textbf{\color{pnpblue}#1}\par\smallskip\ignorespaces}
  {\end{shaded}}

\newcommand{\concept}[1]{\textbf{\color{pnpblue}#1}}
\newcommand{\warn}[1]{\textbf{\color{red!70!black}#1}}

% --- code listings --------------------------------------------------------
\lstset{
  basicstyle=\ttfamily\small,
  breaklines=true,
  columns=fullflexible,
  keepspaces=true,
  showstringspaces=false,
  frame=single,
  framesep=4pt,
  backgroundcolor=\color{calloutbg},
  rulecolor=\color{pnpgray!40},
}

% --- math shortcuts -------------------------------------------------------
\newcommand{\R}{\mathbb{R}}
\newcommand{\diff}{\mathrm{d}}
\newcommand{\eqd}{\overset{!}{=}}
\newcommand{\grad}{\nabla}
\newcommand{\divg}{\nabla\!\cdot}
\newcommand{\eqdef}{\;\coloneqq\;}
\newcommand{\evalat}[2]{\left.#1\right\rvert_{#2}}
\newcommand{\set}[1]{\{#1\}}

% common chemistry / physics symbols
\newcommand{\Hyd}{\mathrm{H}}
\newcommand{\Oxy}{\mathrm{O}}
\newcommand{\HtwoO}{\mathrm{H_2O}}
\newcommand{\HtwoOtwo}{\mathrm{H_2O_2}}
\newcommand{\Otwo}{\mathrm{O_2}}
\newcommand{\Hp}{\mathrm{H^+}}
\newcommand{\OHm}{\mathrm{OH^-}}
\newcommand{\Kp}{\mathrm{K^+}}
\newcommand{\Csp}{\mathrm{Cs^+}}
\newcommand{\SOfour}{\mathrm{SO_4^{2-}}}
\newcommand{\Vrhe}{V_{\mathrm{RHE}}}
\newcommand{\eo}{E^{\circ}}
\newcommand{\Eeq}{E_{\mathrm{eq}}}
\newcommand{\kzero}{k_0}

\newcommand{\onboardingfooter}{%
  \vfill\noindent\rule{\linewidth}{0.4pt}\par
  {\small\color{pnpgray} PNPInverse intern onboarding,
  compiled \today. Source under \texttt{docs/onboarding/}.}\par}
 after \documentclass.

\usepackage[utf8]{inputenc}
\usepackage[T1]{fontenc}
\usepackage{lmodern}
\usepackage[margin=1in]{geometry}
\usepackage{amsmath, amssymb, amsthm}
\usepackage{mathtools}
\usepackage{empheq}
\usepackage{siunitx}
\usepackage{booktabs}
\usepackage{array}
\usepackage{enumitem}
\usepackage{xcolor}
\usepackage{framed}
\usepackage{microtype}
\usepackage{listings}
\usepackage{tikz}
\usetikzlibrary{arrows.meta, positioning, decorations.pathreplacing, calligraphy}
\usepackage[hidelinks]{hyperref}

% --- branding -------------------------------------------------------------
\definecolor{pnpblue}{RGB}{30,70,120}
\definecolor{pnpgray}{RGB}{90,90,90}
\definecolor{calloutbg}{RGB}{245,245,235}
\definecolor{shadecolor}{RGB}{245,245,235}

% --- callouts (uses framed.sty's shaded environment) ----------------------
\newenvironment{callout}[1]
  {\begin{shaded}\noindent\textbf{\color{pnpblue}#1}\par\smallskip\ignorespaces}
  {\end{shaded}}

\newcommand{\concept}[1]{\textbf{\color{pnpblue}#1}}
\newcommand{\warn}[1]{\textbf{\color{red!70!black}#1}}

% --- code listings --------------------------------------------------------
\lstset{
  basicstyle=\ttfamily\small,
  breaklines=true,
  columns=fullflexible,
  keepspaces=true,
  showstringspaces=false,
  frame=single,
  framesep=4pt,
  backgroundcolor=\color{calloutbg},
  rulecolor=\color{pnpgray!40},
}

% --- math shortcuts -------------------------------------------------------
\newcommand{\R}{\mathbb{R}}
\newcommand{\diff}{\mathrm{d}}
\newcommand{\eqd}{\overset{!}{=}}
\newcommand{\grad}{\nabla}
\newcommand{\divg}{\nabla\!\cdot}
\newcommand{\eqdef}{\;\coloneqq\;}
\newcommand{\evalat}[2]{\left.#1\right\rvert_{#2}}
\newcommand{\set}[1]{\{#1\}}

% common chemistry / physics symbols
\newcommand{\Hyd}{\mathrm{H}}
\newcommand{\Oxy}{\mathrm{O}}
\newcommand{\HtwoO}{\mathrm{H_2O}}
\newcommand{\HtwoOtwo}{\mathrm{H_2O_2}}
\newcommand{\Otwo}{\mathrm{O_2}}
\newcommand{\Hp}{\mathrm{H^+}}
\newcommand{\OHm}{\mathrm{OH^-}}
\newcommand{\Kp}{\mathrm{K^+}}
\newcommand{\Csp}{\mathrm{Cs^+}}
\newcommand{\SOfour}{\mathrm{SO_4^{2-}}}
\newcommand{\Vrhe}{V_{\mathrm{RHE}}}
\newcommand{\eo}{E^{\circ}}
\newcommand{\Eeq}{E_{\mathrm{eq}}}
\newcommand{\kzero}{k_0}

\newcommand{\onboardingfooter}{%
  \vfill\noindent\rule{\linewidth}{0.4pt}\par
  {\small\color{pnpgray} PNPInverse intern onboarding,
  compiled \today. Source under \texttt{docs/onboarding/}.}\par}
 after \documentclass.

\usepackage[utf8]{inputenc}
\usepackage[T1]{fontenc}
\usepackage{lmodern}
\usepackage[margin=1in]{geometry}
\usepackage{amsmath, amssymb, amsthm}
\usepackage{mathtools}
\usepackage{empheq}
\usepackage{siunitx}
\usepackage{booktabs}
\usepackage{array}
\usepackage{enumitem}
\usepackage{xcolor}
\usepackage{framed}
\usepackage{microtype}
\usepackage{listings}
\usepackage{tikz}
\usetikzlibrary{arrows.meta, positioning, decorations.pathreplacing, calligraphy}
\usepackage[hidelinks]{hyperref}

% --- branding -------------------------------------------------------------
\definecolor{pnpblue}{RGB}{30,70,120}
\definecolor{pnpgray}{RGB}{90,90,90}
\definecolor{calloutbg}{RGB}{245,245,235}
\definecolor{shadecolor}{RGB}{245,245,235}

% --- callouts (uses framed.sty's shaded environment) ----------------------
\newenvironment{callout}[1]
  {\begin{shaded}\noindent\textbf{\color{pnpblue}#1}\par\smallskip\ignorespaces}
  {\end{shaded}}

\newcommand{\concept}[1]{\textbf{\color{pnpblue}#1}}
\newcommand{\warn}[1]{\textbf{\color{red!70!black}#1}}

% --- code listings --------------------------------------------------------
\lstset{
  basicstyle=\ttfamily\small,
  breaklines=true,
  columns=fullflexible,
  keepspaces=true,
  showstringspaces=false,
  frame=single,
  framesep=4pt,
  backgroundcolor=\color{calloutbg},
  rulecolor=\color{pnpgray!40},
}

% --- math shortcuts -------------------------------------------------------
\newcommand{\R}{\mathbb{R}}
\newcommand{\diff}{\mathrm{d}}
\newcommand{\eqd}{\overset{!}{=}}
\newcommand{\grad}{\nabla}
\newcommand{\divg}{\nabla\!\cdot}
\newcommand{\eqdef}{\;\coloneqq\;}
\newcommand{\evalat}[2]{\left.#1\right\rvert_{#2}}
\newcommand{\set}[1]{\{#1\}}

% common chemistry / physics symbols
\newcommand{\Hyd}{\mathrm{H}}
\newcommand{\Oxy}{\mathrm{O}}
\newcommand{\HtwoO}{\mathrm{H_2O}}
\newcommand{\HtwoOtwo}{\mathrm{H_2O_2}}
\newcommand{\Otwo}{\mathrm{O_2}}
\newcommand{\Hp}{\mathrm{H^+}}
\newcommand{\OHm}{\mathrm{OH^-}}
\newcommand{\Kp}{\mathrm{K^+}}
\newcommand{\Csp}{\mathrm{Cs^+}}
\newcommand{\SOfour}{\mathrm{SO_4^{2-}}}
\newcommand{\Vrhe}{V_{\mathrm{RHE}}}
\newcommand{\eo}{E^{\circ}}
\newcommand{\Eeq}{E_{\mathrm{eq}}}
\newcommand{\kzero}{k_0}

\newcommand{\onboardingfooter}{%
  \vfill\noindent\rule{\linewidth}{0.4pt}\par
  {\small\color{pnpgray} PNPInverse intern onboarding,
  compiled \today. Source under \texttt{docs/onboarding/}.}\par}


\title{\bfseries Electrochemistry Primer for ORR \\
       \large with just enough physical chemistry to read the codebase}
\author{Onboarding doc 01 of 06}
\date{\today}

\begin{document}
\maketitle

\begin{callout}{What this document is for}
The longest of the onboarding docs. It builds, from first
principles, the chemistry vocabulary you need to talk about the
oxygen reduction reaction (ORR), the electric double layer, and the
Butler--Volmer (BV) kinetic equation. Read it once now; you will
return to specific sections as they come up. There are exercises at
the end of each major section --- do them in your notebook before the
meeting; we will go over them together.
\end{callout}

\tableofcontents
\bigskip\hrule\bigskip

% =========================================================================
\section{What an electrochemical cell is}

An \concept{electrochemical cell} is two pieces of metal (or other
electronic conductor) called \concept{electrodes} dipped into a
liquid called the \concept{electrolyte}. The electrolyte is a
solvent (water for us) containing dissolved \concept{ions}: charged
particles that can move through the liquid. The two electrodes are
connected externally through a wire and possibly a power supply.

\textbf{Why is this interesting?} Because \emph{electrons can only
flow in the wire, not in the electrolyte}. So whenever current goes
around the loop, something has to convert electron flow at the
electrode--electrolyte interface into ion flow in the electrolyte.
That ``something'' is a chemical reaction at the interface. Different
reactions happen at different rates and at different applied
voltages, and that is the entire content of electrochemistry.

\subsection{Half-cells, oxidation, reduction}

Each electrode hosts a \concept{half-reaction}. By convention,
\concept{reduction} (gain of electrons) happens at the
\concept{cathode}; \concept{oxidation} (loss of electrons) happens
at the \concept{anode}. The classical mnemonic is
``\textsc{red cat}, \textsc{an ox}.'' For the ORR cell we model:

\begin{itemize}[leftmargin=*]
\item \textbf{Cathode (the disk):} oxygen is reduced.
  $\Otwo + n\Hp + n e^- \rightarrow \text{products}$.
\item \textbf{Anode (somewhere far away in the cell, often a Pt
  counter electrode):} water is oxidized to $\Otwo$, the so-called
  oxygen evolution reaction.
\end{itemize}

We model only the cathode region because that is where the
interesting physics is. The counter-electrode is treated as a
boundary condition implicitly through the bulk concentrations.

\subsection{Faraday's law (the dimensional anchor)}

The central conversion factor between electricity and chemistry is
\concept{Faraday's constant}:
\[
  F \approx 96\,485 \;\; \mathrm{C/mol\;e^-}.
\]
Faraday's law of electrolysis states that the amount of substance
reacted is proportional to the charge passed:
\[
  \text{(moles reacted)} = \frac{Q}{n_e\,F},
\]
where $Q$ is the integrated current (in coulombs) and $n_e$ is the
number of electrons transferred per molecule. Differentiating:
\begin{equation}
  i = n_e\,F\,r_{\text{rxn}},
  \label{eq:faraday}
\end{equation}
where $i$ is the current density (A/m$^2$) and $r_{\text{rxn}}$ is
the molar reaction rate per unit area (mol/m$^2$/s). This equation
is the bridge between the chemical models in our code and the
\textit{measured} quantity $i$.

\subsection{Convention: cathodic current is negative}

Electrochemistry sign conventions are a constant source of bugs.
The convention this codebase uses (and the one most ORR papers use):

\begin{callout}{Sign convention}
\begin{itemize}[leftmargin=*]
\item Current density $i > 0$ for \textbf{anodic} (oxidative) flow.
\item $i < 0$ for \textbf{cathodic} (reductive) flow --- ORR makes
      $i < 0$.
\item The disk-current curves you will see plotted are typically
      plotted as $|i|$ vs.\ $V$, so visually the cathodic ``ramp up''
      goes \textit{up}, even though the underlying $i$ is going more
      negative.
\end{itemize}
\end{callout}

\section{Equilibrium electrochemistry}

\subsection{Standard reduction potential $\eo$}

Each half-reaction has a thermodynamic equilibrium voltage. By
convention, half-reactions are written as reductions and tabulated
relative to the \concept{standard hydrogen electrode (SHE)}, which
is defined to have $\eo \equiv 0$. For the ORR we care about two:

\begin{align}
  \Otwo + 4\Hp + 4 e^- &\rightleftharpoons 2\,\HtwoO &
  &\eo_{4e} = +1.23\;\Vrhe \\
  \Otwo + 2\Hp + 2 e^- &\rightleftharpoons \HtwoOtwo &
  &\eo_{2e} = +0.695\;\Vrhe
\end{align}

These are the values from Ruggiero 2022 \S 1, reported on the
\emph{reversible hydrogen electrode} (RHE) scale --- the convention
used throughout this codebase, the deck, and the ORR literature. The
RHE scale absorbs the pH dependence of any proton-coupled electron
transfer (PCET) reaction. Numerically, at pH $= 0$ these coincide
with the standard reduction potentials vs SHE; at other pH the
SHE-referenced values shift by $-0.059\,\mathrm{V}\cdot\mathrm{pH}$,
exactly cancelling the RHE-vs-SHE shift, so the values quoted above
are pH-independent. The next subsection makes the SHE/RHE conversion
explicit.

The 4-electron pathway is more thermodynamically favorable (higher
$\eo$) but also requires breaking the strong $\Oxy\text{--}\Oxy$
bond, which can be slow --- hence the kinetic ``why does the 2e
pathway happen at all?'' question that drives the field.

\subsection{The Nernst equation}

If the reactants and products are not at standard concentration
(1 mol/L for solutes, 1 atm for gases), the equilibrium voltage
shifts. For a generic half-reaction
$\nu_O \,\mathrm{O} + n_e e^- \rightleftharpoons \nu_R\,\mathrm{R}$,
the \concept{Nernst equation} gives:
\begin{equation}
  \Eeq = \eo - \frac{RT}{n_e F} \ln\!\left(
    \frac{a_R^{\nu_R}}{a_O^{\nu_O}}\right),
  \label{eq:nernst}
\end{equation}
where $a_i$ is the \concept{activity} of species $i$ (close to
concentration in dilute solution, with corrections we will discuss).
The prefactor $RT/F$ at room temperature is the
\concept{thermal voltage}:
\[
  V_T \;\eqdef\; \frac{RT}{F} \approx 25.7 \;\text{mV \, at \,} T = 298\,\mathrm{K}.
\]
You will see $V_T$ everywhere in the code (\texttt{V\_T} in
\texttt{scripts/\_bv\_common.py}).

\subsection{Reference electrodes: SHE, Ag/AgCl, RHE}

In the lab, nobody uses a hydrogen electrode --- it requires
bubbling explosive H$_2$ at exactly 1 atm over a Pt black surface.
Real experiments use practical references like
\concept{Ag/AgCl in saturated KCl}, then \emph{convert} to a
common pH-aware reference, the \concept{reversible hydrogen
electrode (RHE)}:
\[
  V_{\mathrm{RHE}} = V_{\mathrm{Ag/AgCl}}
                   + 0.197\,\mathrm{V}
                   + 0.059\,\mathrm{V}\cdot \mathrm{pH}.
\]
The 0.059 V/pH conversion is just $V_T \ln 10 \approx 59$ mV at room
temperature, applied to the Nernstian pH dependence of the H$_2$
half-reaction. \textbf{Every voltage in our code and in the deck data
is reported as $\Vrhe$.} If you ever see a paper quoting voltages
``vs.\ SHE'' or ``vs.\ Ag/AgCl,'' convert before comparing.

\subsection{Overpotential $\eta$}

When the applied voltage $V$ at an electrode differs from the
local equilibrium $\Eeq$, a current flows. The driving force is the
\concept{overpotential}:
\begin{equation}
  \eta \;\eqdef\; V - \Eeq.
  \label{eq:overpotential}
\end{equation}

\textbf{Note the sign:} $\eta < 0$ drives the reduction direction
(cathodic), $\eta > 0$ drives oxidation (anodic). Most of the ORR
voltage range we care about has $\eta < 0$.

\paragraph{Exercise 2.1.}
For the 4e ORR, the standard reduction potential at pH 0 is
$E^{\circ}(\mathrm{SHE}) = +1.229$ V. Apply the Nernst equation at
bulk pH 4 ($[\Hp] = 10^{-4}$ M, $[\Otwo] = 1.2 \times 10^{-3}$ M,
$P_{\Otwo} = 1$ atm; activity of liquid water $= 1$) to compute
$\Eeq$ vs.\ SHE. Then convert to $\Vrhe$ using the RHE/SHE relation
in \S 2.3. You should recover $\Eeq \approx +1.23\;\Vrhe$,
confirming that the RHE-referenced potential of any PCET reaction
with $n_{\Hp}/n_e = 1$ is pH-independent.

% =========================================================================
\section{The electric double layer (EDL)}

This section is the bridge between equilibrium thermodynamics and
the much stranger world of the electrode-electrolyte
\textit{interface}.

\subsection{Why an interface develops a charge layer}

When a metal electrode is immersed in an electrolyte and held at
some applied potential $V$, electrons accumulate (or deplete) at the
metal surface. To maintain overall electroneutrality on the
solution side within a few nanometers, ions of the opposite sign
crowd toward the surface and ions of the same sign are repelled.
The result is a thin region of net charge in the electrolyte adjacent
to the metal, balancing the surface electronic charge: this is the
\concept{electric double layer (EDL)}.

The EDL controls almost everything we care about:
\begin{itemize}[leftmargin=*]
\item It sets the local concentrations of $\Hp$, cations, etc.\ at
      the electrode surface --- where the chemistry happens.
\item It absorbs a portion of the applied voltage, so the
      ``available driving force'' for the reaction is less than
      $V - \Eeq^{\mathrm{bulk}}$.
\item It depends on cation identity and salt concentration, which is
      one mechanism for the cation effect.
\end{itemize}

\subsection{Three increasingly realistic EDL models}

\paragraph{Helmholtz (1853).}
A single rigid sheet of counter-ions sits one ionic radius from the
metal, like a parallel-plate capacitor. Predicts a constant
capacitance independent of voltage and concentration. Wrong, but
simple.

\paragraph{Gouy--Chapman (1910--13).}
The counter-ion ``sheet'' is allowed to spread out as a continuous
\concept{diffuse layer} obeying a Boltzmann distribution:
\begin{equation}
  c_i(x) = c_i^{\mathrm{bulk}}\,\exp\!\left(-\frac{z_i F \phi(x)}{RT}\right),
  \label{eq:boltzmann}
\end{equation}
where $\phi(x)$ is the local electrostatic potential measured
relative to the bulk and $z_i$ is the ion charge number. The width
of this layer is the \concept{Debye length}:
\begin{equation}
  \lambda_D = \sqrt{\frac{\varepsilon RT}{F^2 \sum_i z_i^2 c_i^{\mathrm{bulk}}}},
  \label{eq:debye}
\end{equation}
which for our 0.3 M ionic strength electrolyte gives
$\lambda_D \approx 0.55$ nm. Submerged in this region, a continuum
description still works (it has been re-derived from first
principles; see Bard \& Faulkner Ch.\ 13).

\paragraph{Stern (1924).}
Combines the two: a rigid \concept{compact layer} (Stern layer) of
finite thickness right against the electrode, beyond which the
Gouy--Chapman diffuse layer takes over. The Stern layer is treated
as a parallel-plate capacitor with capacitance $C_S$
(F/m$^2$). Most of the applied voltage drops across the Stern
layer at high overpotential.

\paragraph{Bikerman (1942).}
Adds the missing physics that ions have finite size: the
exponential Boltzmann factor (\ref{eq:boltzmann}) cannot pile up
arbitrary amounts of counter-ion at the surface, because once you
hit close-packing the local steric pressure resists further
crowding. Bikerman's modification:
\begin{equation}
  c_i(x) = \frac{c_i^{\mathrm{bulk}}\,e^{-z_i F\phi(x)/RT}}
                {1 + \sum_j \frac{c_j^{\mathrm{bulk}}}{c_{\max}}
                       \left(e^{-z_j F\phi(x)/RT} - 1\right)}.
  \label{eq:bikerman}
\end{equation}
At large $|\phi|$, this saturates at $c_i \to c_{\max}$ instead of
diverging. The codebase uses a multi-species Bikerman closure
defined in \texttt{Forward/bv\_solver/boltzmann.py}; we will look at
it in doc 03.

\subsection{A picture of the EDL}

\begin{center}
\begin{tikzpicture}[
  every node/.style={font=\small},
  ion/.style={circle, draw, minimum size=4mm, inner sep=0pt},
]
% metal electrode
\fill[gray!50] (0,-1.2) rectangle (0.6,1.2);
\node[rotate=90] at (0.3,0) {\bfseries metal};

% Stern layer
\draw[dashed] (1.4,-1.2) -- (1.4,1.2);
\draw[dashed] (2.6,-1.2) -- (2.6,1.2);
\node[align=center] at (2.0,1.5) {Stern \\ (compact)};

% diffuse layer
\node[align=center] at (4.5,1.5) {Diffuse \\ (Gouy--Chapman)};
\draw[dashed] (6.5,-1.2) -- (6.5,1.2);
\node[align=center] at (8.2,1.5) {Bulk \\ electrolyte};

% counter-ions in Stern + diffuse
\node[ion, fill=blue!30] at (2.0, 0.7) {$+$};
\node[ion, fill=blue!30] at (2.0, 0.0) {$+$};
\node[ion, fill=blue!30] at (2.0,-0.7) {$+$};
\node[ion, fill=blue!30] at (3.2, 0.4) {$+$};
\node[ion, fill=blue!30] at (3.2,-0.4) {$+$};
\node[ion, fill=blue!30] at (4.5, 0.5) {$+$};
\node[ion, fill=blue!30] at (4.5,-0.6) {$+$};
\node[ion, fill=red!40]  at (5.7, 0.6) {$-$};
\node[ion, fill=blue!30] at (5.5,-0.4) {$+$};
\node[ion, fill=red!40]  at (5.9,-0.7) {$-$};
\node[ion, fill=blue!30] at (7.2, 0.5) {$+$};
\node[ion, fill=red!40]  at (7.5,-0.4) {$-$};
\node[ion, fill=blue!30] at (8.0, 0.0) {$+$};
\node[ion, fill=red!40]  at (8.4, 0.7) {$-$};

% potential profile
\draw[thick,blue!70!black] plot[smooth, tension=0.6]
   coordinates {(0.6,1.0) (1.4,0.7) (2.6,0.5) (4.0,0.18) (6.5,0.04) (8.5,0.0)};
\node[blue!70!black] at (4,-1.6) {$\phi(x)$ falls off over $\lambda_D \sim 0.5$--$1$ nm};

\draw[<-, thick] (0.6,-1.4) -- (0.6,-1.7);
\node[below] at (0.6,-1.7) {$x=0$ (metal surface)};
\end{tikzpicture}
\end{center}

The metal sits at $x=0$ at the left. The Stern layer is one or two
ionic radii thick (about 0.3 nm). Beyond it the diffuse layer
extends a few Debye lengths into the electrolyte before the
potential and concentrations relax to their bulk values.

\paragraph{Exercise 3.1.} For $z_i = +1$, $V_T = 25.7$ mV at room
temperature, what is the ratio $c_i(0)/c_i^{\mathrm{bulk}}$ from the
plain Boltzmann form (\ref{eq:boltzmann}) at $\phi(0) = -0.3$ V? Why
is this physically absurd, and what does Bikerman do about it?

% =========================================================================
\section{Electrochemical potential}

\subsection{Definition}

When a charged species $i$ sits in a region of nonzero electric
potential $\phi$, its true thermodynamic ``willingness to be there''
is not just its chemical potential $\mu_i$ but the
\concept{electrochemical potential}:
\begin{equation}
  \tilde\mu_i = \mu_i^{\circ} + RT \ln a_i + z_i F \phi.
  \label{eq:echempot}
\end{equation}
\textit{This is the central thermodynamic quantity of the entire
field}, and it is also the variable our solver uses internally for
$\Hp$ (the so-called \texttt{logc\_muh} formulation: ``log-c with
$\mu_H$''). We will see in doc 02 why that choice makes Newton's
method behave better.

At equilibrium, $\tilde\mu_i$ is uniform across the electrolyte. So
if you know $\phi(x)$, you can solve for $c_i(x)$ from
$\tilde\mu_i = \tilde\mu_i^{\mathrm{bulk}}$, which immediately
recovers the Boltzmann distribution (\ref{eq:boltzmann}).

\subsection{The thermal voltage $V_T$ as a unit}

Notice that everywhere $z F \phi$ appears, it is divided by $RT$ to
form $z\phi/V_T$. So in our nondimensionalized code, voltages are
always reported in units of $V_T$, and the variable
\texttt{phi\_hat} $= \phi/V_T$ is what the residual sees. Doc 02 has
the full nondimensionalization.

% =========================================================================
\section{Butler--Volmer kinetics}

\subsection{The classical form}

For a single elementary electron-transfer reaction
\[
  \mathrm{O} + n_e e^- \rightleftharpoons \mathrm{R},
\]
\concept{transition-state theory} (sketch in any electrochemistry
textbook, e.g.\ Bard \& Faulkner Ch.\ 3) gives the net forward
current density:
\begin{equation}
  i = i_0 \left[
       \exp\!\left(\frac{(1-\alpha) n_e F \eta}{RT}\right)
     - \exp\!\left(\frac{-\alpha n_e F \eta}{RT}\right)
  \right],
  \label{eq:bv}
\end{equation}
where:
\begin{itemize}[leftmargin=*]
\item $\eta = V - \Eeq$ is the overpotential.
\item $\alpha \in (0,1)$ is the \concept{symmetry factor} or
      \concept{transfer coefficient} (often $\alpha \approx 0.5$).
\item $i_0$ is the \concept{exchange current density} at $\eta = 0$;
      it depends on $\kzero$ (the rate constant) and the local
      concentrations.
\end{itemize}

We use a slightly different decomposition that scales better
numerically (the ``log-rate'' form, \texttt{bv\_log\_rate=True}):
\begin{equation}
  \log r = \log \kzero + \sum_p p \cdot \log a_p
           - \alpha n_e \, \frac{\eta}{V_T},
  \label{eq:logbv}
\end{equation}
where $p$ runs over the species in the rate law with stoichiometric
coefficient $\nu_p$. The reason for log-rate is that $\kzero$ varies
over many orders of magnitude in our calibration, and Newton iterates
on $\log r$ stay smooth even in regions where $r$ itself underflows
double precision.

\subsection{Tafel limits}

When $|\eta|$ is large enough that one of the exponentials in
(\ref{eq:bv}) dwarfs the other, the equation reduces to the
\concept{Tafel form}:
\[
  \log_{10}|i| \approx \log_{10} i_0 + \frac{(1-\alpha) n_e}{2.303\,V_T}\,\eta
  \quad\text{(anodic)},
\]
giving a straight line of slope $\sim 60/[(1-\alpha) n_e]$ mV/decade.
Plotting $\log|i|$ vs.\ $\eta$ and fitting straight-line slopes is
how experimentalists historically measured $\alpha$ and $\kzero$.
Look up \concept{Tafel slope} when you see it in commit messages.

\subsection{Multi-step ORR and the parallel topology}

The ORR is not a single elementary step. It proceeds through
several adsorbed intermediates:
\[
  \Otwo \to {}^*\!\Otwo \to {}^*\!\mathrm{OOH} \to {}^*\!\Oxy
  \to {}^*\!\mathrm{OH} \to \HtwoO,
\]
where ${}^*$ denotes a surface-adsorbed species. At any point along
this chain the system can branch off: if the $*\!\mathrm{OOH}$
intermediate desorbs as $\HtwoOtwo$, you exit early with the
2-electron product; if it stays bound and gets further reduced, you
proceed to water through the 4-electron path. Following Ruggiero
2022 \S 1 we do not resolve the intermediates; instead we lump them
into two \emph{parallel} effective Butler--Volmer reactions,
\textbf{R$_{\text{2e}}$} and \textbf{R$_{\text{4e}}$}, each with
its own $\eo$, $\kzero$, and $\alpha$:
\begin{align}
  \mathrm{R_{2e}}: && \Otwo + 2\Hp + 2 e^- &\to \HtwoOtwo & \eo &= 0.695\;\Vrhe \\
  \mathrm{R_{4e}}: && \Otwo + 4\Hp + 4 e^- &\to 2\,\HtwoO & \eo &= 1.23\;\Vrhe
\end{align}
Both reactions consume $\Otwo$. They share no surface intermediate
in the lumped model. The total disk current is
$i_D = 2F r_{2e} + 4F r_{4e}$. The peroxide that escapes is
$2F r_{2e}$.

\warn{Historical hazard:} the codebase used to use a different,
\emph{sequential} reaction set (R$_0$: $\Otwo \to \HtwoOtwo$;
R$_1$: $\HtwoOtwo \to \HtwoO$). That set was wrong --- it
conflicts with Ruggiero 2022 --- and was retired in May 2026
(commit M3a.2). You will still see ``R$_0$/R$_1$'' in old commits,
old plot legends, and old constants like \texttt{E\_eq\_r1=0.68}.
\textbf{All deck-aligned work uses the parallel 2e/4e set.}

\subsection{Selectivity and electron number}

In an RRDE experiment (next section) you measure two currents: the
disk current $I_D$ (negative, all the cathode reduction at the disk)
and the ring current $I_R$ (positive, peroxide that escaped is being
re-oxidized at the downstream ring). From these:
\begin{align}
  \mathrm{H_2O_2}\,\% &= \frac{200 \cdot (I_R / N)}{|I_D| + I_R/N} \label{eq:selectivity}\\
  n_e &= \frac{4 \cdot |I_D|}{|I_D| + I_R/N}
\end{align}
where $N \in (0,1)$ is the geometry-determined
\concept{collection efficiency}: only a fraction of the peroxide
made at the disk reaches the ring before being washed downstream.
For our deck, $N = 0.224$.

If \emph{all} of the disk current came through R$_{2e}$,
selectivity is 100\% and $n_e = 2$. If \emph{all} comes through
R$_{4e}$, selectivity is 0\% and $n_e = 4$. Real experiments fall
in between.

\paragraph{Exercise 5.1.} Derive the factor of 200 in
(\ref{eq:selectivity}) by counting electrons. Hint: you have to
correct $I_R$ by $1/N$ to recover the actual $\HtwoOtwo$ flux out
of the disk; then the ratio (peroxide-equivalent current at the
disk) / (total disk current as 4e-equivalent) gives the fraction.

% =========================================================================
\section{Mass transport: how reactants reach the electrode}

The BV equation tells you the rate at which the surface chemistry
\textit{would} consume reactants if there were unlimited supply. In
reality, $\Otwo$ has to diffuse from the bulk to the electrode and
$\Hp$ has to as well. Above some current density, you run out of
reactant at the surface and the kinetic equation is no longer the
bottleneck --- you are \concept{mass-transport limited}.

\subsection{The flux equation: Nernst--Planck}

For a dilute species $i$, the flux per unit area is the sum of three
contributions:
\begin{equation}
  \mathbf{J}_i = \underbrace{-D_i \grad c_i}_{\text{diffusion}}
                 \underbrace{- \frac{z_i F}{RT}D_i c_i \grad\phi}_{\text{migration}}
                 + \underbrace{c_i \mathbf{v}}_{\text{convection}}.
  \label{eq:np}
\end{equation}
Mass conservation gives $\partial_t c_i + \divg \mathbf{J}_i = 0$.
This is the \concept{Nernst--Planck equation}, the ``NP'' in PNP.
We almost always work in steady state ($\partial_t = 0$) and 1D, so
this becomes a coupled boundary-value problem.

\subsection{The rotating disk electrode (RDE)}

The disk in our experiment rotates at $\omega = 1600$ rpm. This
creates a well-defined hydrodynamic flow that brings fresh
electrolyte from the bulk toward the disk and then flings it
radially outward. Levich (1962) showed that the convection and the
diffusion combine into an effective \concept{diffusion layer
thickness}:
\begin{equation}
  \delta \approx 1.612\, D_i^{1/3} \nu^{1/6} \omega^{-1/2},
  \label{eq:levich-thickness}
\end{equation}
where $\nu$ is the kinematic viscosity of the electrolyte. For
$\Otwo$ in our system, $\delta \approx 16$ $\mu$m at 1600 rpm.

The corresponding maximum (Levich-limited) current density when the
surface reactant is fully depleted is:
\begin{equation}
  i_{\mathrm{lim}} = 0.62 \cdot n_e F \, D_i^{2/3}\nu^{-1/6}\omega^{1/2}\, c_i^{\mathrm{bulk}}.
  \label{eq:levich}
\end{equation}
For our $\Otwo$ system at 1600 rpm,
$i_{\mathrm{lim}}^{(4e)} \approx 5.7$ mA/cm$^2$. Whenever you see a
plateau in a current-voltage curve, the first thing to check is
whether you have hit Levich.

\subsection{Why we use a 1D slab geometry in the model}

The RDE setup creates flow that is ``\concept{uniformly accessible}''
across the disk: every point on the disk surface sees the same
$\delta$ and the same bulk feed. So we can model the whole disk
problem as a 1D slab from $x=0$ (electrode) to $x=L_{\mathrm{eff}}$
(bulk), with bulk concentrations as Dirichlet boundary conditions at
$x = L_{\mathrm{eff}}$. The convection term in (\ref{eq:np}) is
absorbed by choosing $L_{\mathrm{eff}}$ to match the Levich
$\delta$. Our default is
$L_{\mathrm{eff}} \approx 16$--$100\,\mu$m depending on the study.

\subsection{The ring of the RRDE}

Add a second concentric Pt ring electrode just downstream of the
disk, held at a much more positive voltage (+1.2 V vs.\ RHE) where
\emph{any} $\HtwoOtwo$ that reaches it gets oxidized back to
$\Otwo$. The ring current $I_R$ then directly measures the flux of
peroxide leaving the disk. The fraction of disk-released peroxide
that reaches the ring is the \emph{collection efficiency} $N$ you saw
above (eq.~\ref{eq:selectivity}). Geometry-only; calibrated with a
known redox couple.

\begin{center}
\begin{tikzpicture}[scale=0.9]
\fill[gray!30] (-3.5,0) rectangle (3.5,1.5);
\fill[blue!50] (-3.5,0) rectangle (-2.0,-0.4);
\fill[orange!80] (-1.0,0) rectangle (1.0,-0.4);
\fill[blue!50] (2.0,0) rectangle (3.5,-0.4);
\node at (0,-0.7) {\small disk (carbon, ORR)};
\node at (-2.7,-0.7) {\small ring};
\node at ( 2.7,-0.7) {\small ring};
\draw[->, thick, blue!70!black] (0,1.2) -- (0,0.2);
\node[blue!70!black, right] at (0.1,0.7) {bulk feed (rotation)};
\draw[->, thick, red!70!black] (0.5,0.2) .. controls (1.6,0.6) .. (2.6,0.2);
\node[red!70!black, above] at (1.4,0.5) {\scriptsize peroxide flux $\to$ ring};
\node at (-3.5,1.0) [right] {electrolyte};
\end{tikzpicture}
\end{center}

\paragraph{Exercise 6.1.} If $|I_D| = 2.5$ mA, $I_R = 0.090$ mA,
$N = 0.224$: compute selectivity $\HtwoOtwo$\% and effective
electron number $n_e$.

% =========================================================================
\section{Cation effects: the Phase 6$\beta$ story}

This section is the bridge to current research. It is the
\textit{specific} chemistry our group is working on, and you will
see all of these terms in commits, plots, and the
\texttt{docs/phase6/} directory.

\subsection{The empirical observation}

Hold everything fixed --- catalyst, pH, total ionic strength, applied
voltage --- and only change which alkali cation is in the salt:
$\mathrm{Li_2SO_4} \to \mathrm{Na_2SO_4} \to \mathrm{K_2SO_4} \to
\mathrm{Cs_2SO_4}$. The disk current and the peroxide selectivity
shift substantially. This is the \concept{cation effect}. It has
been reproducibly observed for decades and across many catalysts. The
mystery: cations don't appear in the ORR half-reaction.

\subsection{The currently leading hypothesis: cation hydrolysis at
                                              the OHP}

The hypothesis our group works under (following Singh 2016, Co \&
Zhang 2019, Ringe 2019) is that the ``inert'' cation participates in
a side reaction at the outer edge of the Stern layer (the
\concept{outer Helmholtz plane}, OHP). The hydrated cation
$\mathrm{M(H_2O)_n^+}$ can give up a proton in a hydrolysis
equilibrium:
\begin{equation}
  \mathrm{M(H_2O)_n^+} + \HtwoO
  \;\rightleftharpoons\;
  \mathrm{M(OH)(H_2O)_{n-1}^{0}} + \mathrm{H_3 O^+}.
  \label{eq:catnion-hyd}
\end{equation}
Two things make this matter near a polarized cathode:

\begin{enumerate}[leftmargin=*]
\item \textbf{Field-dependent pK$_a$.} Singh 2016 derived that the
      bulk pK$_a$ of (\ref{eq:catnion-hyd}) shifts by an
      electric-field-dependent amount near the cathode:
      \begin{equation}
        \mathrm{pKa}(\sigma) = \mathrm{pKa}_{\mathrm{bulk}}
          + 2 A z \sigma\, r_{H\!-\!El}
            \left(1 - \frac{r_{M-O}^2}{r_{H\!-\!El}^2}\right),
        \label{eq:singh}
      \end{equation}
      with $\sigma$ the surface charge density at the electrode,
      $r_{M-O}$ the cation--O distance, $r_{H\!-\!El}$ the
      hydration-shell-H-to-electrode distance, and $A = 620.32$ pm.
      The geometry makes $\Delta\mathrm{pKa} < 0$ for cathodic
      $\sigma > 0$: the hydrolysis becomes much more favorable near
      the surface than in the bulk. For Cs$^+$ near a Cu cathode,
      pK$_a$ drops from $14.8$ in bulk to $\sim 4.3$ at typical
      ORR field strengths.
\item \textbf{Cation identity sets the bulk pK$_a$.} Smaller, more
      strongly hydrating cations (Li$^+$) have lower bulk pK$_a$
      values; larger ones (Cs$^+$) have higher ones (see Singh
      Table S1). At a fixed bulk pH and fixed $\sigma$, this
      changes \emph{which} cations are actively releasing protons
      at the OHP and which are not.
\end{enumerate}

The protons released by (\ref{eq:catnion-hyd}) are released exactly
where the ORR cathodic chemistry is consuming them, locally
shielding the surface from the otherwise severe alkalinization
($\HtwoO + e^- \to \tfrac{1}{2}\Otwo + \OHm$ raises local pH a lot).
This converts the cation from a passive spectator into an active,
local pH-buffering species, even though it never enters the formal
ORR equation.

\subsection{Why this is hard to model}

The cation enters through:
\begin{itemize}[leftmargin=*]
\item the EDL physics (Bikerman steric crowding at the OHP),
\item a coupled \concept{coverage} variable
      $\Gamma_{\mathrm{MOH}}$ on the OHP that obeys
      its own ODE in time (rate of (\ref{eq:catnion-hyd}) forward
      vs.\ reverse, both functions of the local field, which is
      itself a function of $\Gamma$ via charge balance, etc.),
\item a feedback into the proton residual at the boundary that
      changes the BV reaction rates.
\end{itemize}
This nonlinear feedback loop is what \concept{Phase 6$\beta$ v9},
\concept{v10a}, and \concept{v10b} have been progressively wiring
into the solver. Doc 04 walks the chronology; for now, the bullet
to remember is: \emph{the next big science our model needs to
reproduce is cation-dependent peroxide selectivity, and the
mechanism we are betting on is field-dependent hydrolysis of
hydrated alkali cations at the OHP.}

\paragraph{Exercise 7.1.} Given Singh's bulk pK$_a$ values
\{Li: 13.6, Na: 14.2, K: 14.5, Cs: 14.8\} and a uniform
$\Delta \mathrm{pKa} = -10$ shift near a polarized Cu cathode, which
of the four cations is most strongly buffering at $\Vrhe = -0.2$ V
in a pH 4 bulk? (Hint: ask yourself which cations have a local
pK$_a$ near or below the local pH, since only those will donate
appreciable proton.)

% =========================================================================
\section{Glossary of terms you will see in commits}

\begin{center}
\begin{tabular}{p{0.20\linewidth} p{0.74\linewidth}}
\toprule
\textbf{Term} & \textbf{Meaning} \\
\midrule
ORR     & Oxygen reduction reaction. $\Otwo \to \HtwoO$ or $\HtwoOtwo$. \\
RRDE    & Rotating ring--disk electrode --- the experimental setup. \\
$\Vrhe$ & Voltage vs.\ reversible hydrogen electrode. \\
$\eta$  & Overpotential, $V - \Eeq$. Negative is cathodic. \\
$\kzero$ & Reaction-rate prefactor (intrinsic activity). \\
$\alpha$ & Symmetry factor in BV, typically $\sim 0.5$. \\
$E^{\circ}$ & Standard reduction potential. \\
EDL     & Electric double layer. \\
OHP     & Outer Helmholtz plane: the boundary between the Stern and diffuse layers. \\
$C_S$   & Stern (compact-layer) capacitance, F/m$^2$. We use 0.20 F/m$^2$. \\
$\lambda_D$ & Debye length: width of the diffuse layer, $\sim 0.55$ nm here. \\
$\delta$    & Levich diffusion-layer thickness, $\sim 16$ $\mu$m at 1600 rpm. \\
$L_{\mathrm{eff}}$ & Slab thickness in our 1D model, set to $\delta$. \\
$N$         & RRDE collection efficiency, $0.224$ here. \\
$n_e$       & Effective electron number, between 2 and 4. \\
PNP    & Poisson--Nernst--Planck. The PDE system. \\
BV     & Butler--Volmer. The kinetic boundary condition. \\
Bikerman & Steric correction to the Boltzmann distribution. \\
Phase 6$\alpha$ & Project epoch: water self-ionization landed. \\
Phase 6$\beta$  & Project epoch: cation hydrolysis at OHP (\textit{ongoing}). \\
\bottomrule
\end{tabular}
\end{center}

\section{Where to read more}

\begin{itemize}[leftmargin=*]
\item Bard \& Faulkner, \textit{Electrochemical Methods} (2nd ed.):
      Chapters 1, 3, 12, 13. The standard textbook; chapters 1 and 3
      give you everything in this primer in proper depth.
\item Ruggiero 2022 PDF
      (\texttt{docs/papers/Ruggiero2022\_JCatal\_manuscript.pdf}):
      our experimental anchor. Read \S 1, \S 2, and Fig.\ 1B.
\item Internal cheat-sheet
      (\texttt{docs/papers/Ruggiero2022\_JCatal\_source\_paper.md}):
      compact list of every number in the paper we use.
\item Singh 2016 SI extraction
      (\texttt{docs/phase6/singh\_2016\_pka\_formula.md}): how
      Eq.~\ref{eq:singh} is derived and verified against
      experimental titration data for the alkali metals.
\item CMK-3 capacitance literature
      (\texttt{docs/phase6/CMK3\_capacitance\_literature.md}): the
      citation chain behind our $C_S = 0.20$ F/m$^2$ choice.
\end{itemize}

\onboardingfooter
\end{document}
